\clearpage
\fancyhead[L]{\small INHALT.}
\SourceLeaf{660}
\begingroup
\raggedright
\begin{center}\textbf{INHALT.}\par\medskip
\textbf{GAUSS WERKE BAND VI. ASTRONOMISCHE ABHANDLUNGEN.}\end{center}
\par\medskip
\noindent\textit{Abhandlungen.}\par\smallskip\nopagebreak[4]
\noindent\hangindent=1em Disquisitio de elementis ellipticis Palladis ex oppositionibus annorum 1803, 1804, 1805, 1807, 1808, 1809 \mbox{.\ .\ .\ .\ .\ .\ .\ .\ .\ .\ .\ .\ .\ .\ .\ .} \hfill\mbox{1810 Novbr. 25}\quad\mbox{Seite 1}\par\smallskip
\noindent\hangindent=1em Observationes cometae secundi a. 1813 in observatorio Gottingensi factae, adjectis nonnullis adnotationibus circa calculum orbitarum parabolicarum \hfill\mbox{1813 Sept. 10}\quad\mbox{--- 25}\par\smallskip
\noindent\hangindent=1em Methodus peculiaris elevationem poli determinandi \mbox{.\ .\ .\ .\ .\ .\ .\ .\ .} \hfill\mbox{1808}\quad\mbox{--- 37}\par\smallskip
\par\medskip
\noindent\textit{Anzeigen eigner Schriften.}\par\smallskip\nopagebreak[4]
\noindent\hangindent=1em Theoria motus corporum coelestium in sectionibus conicis Solem ambientium \hfill\mbox{1809 Juni 17}\quad\mbox{--- 53}\par\smallskip
\noindent\hangindent=1em Disquisitio de elementis ellipticis Palladis ex oppositionibus annorum 1803, 1804, 1805, 1807, 1808, 1809 \mbox{.\ .\ .\ .\ .\ .\ .\ .\ .\ .\ .\ .\ .\ .\ .\ .} \hfill\mbox{1810 Dec. 10}\quad\mbox{--- 61}\par\smallskip
\noindent\hangindent=1em Observationes cometae secundi a. 1813 in observatorio Gottingensi factae, adjectis nonnullis adnotationibus circa calculum orbitarum parabolicarum \hfill\mbox{1814 Jan. 3}\quad\mbox{--- 65}\par\smallskip
\noindent\hangindent=1em Methodus peculiaris elevationem poli determinandi \mbox{.\ .\ .\ .\ .\ .\ .\ .\ .} \hfill\mbox{1808 Dec. 5}\quad\mbox{--- 68}\par\smallskip
\par\medskip
\noindent\textit{Verschiedene Aufsätze über Astronomie.}\par\smallskip\nopagebreak[4]
\noindent\hangindent=1em Berechnung des Osterfestes \mbox{.\ .\ .\ .\ .\ .\ .\ .\ .\ .\ .\ .\ .\ .\ .\ .\ .\ .} \hfill\mbox{1800 Aug.}\quad\mbox{--- 73}\par\smallskip
\noindent\hangindent=1em Berechnung des jüdischen Osterfestes \mbox{.\ .\ .\ .\ .\ .\ .\ .\ .\ .\ .\ .\ .\ .} \hfill\mbox{1802 Mai}\quad\mbox{--- 80}\par\smallskip
\noindent\hangindent=1em Noch etwas über die Bestimmung des Osterfestes \mbox{.\ .\ .\ .\ .\ .\ .\ .\ .\ .} \hfill\mbox{1807 Sept. 12}\quad\mbox{--- 82}\par\smallskip
\noindent\hangindent=1em Vorschriften, um aus der geocentrischen Länge und Breite eines Himmelskörpers, dem Orte seines Knotens, der Neigung der Bahn, der Länge der Sonne und ihrem Abstande von der Erde abzuleiten: Des Himmelskörpers heliocentrische Länge in der Bahn, wahren Abstand von der Sonne und wahren Abstand von der Erde \mbox{.\ .\ .\ .\ .\ .\ .\ .\ .\ .\ .} \hfill\mbox{1802 Juni}\quad\mbox{--- 87}\par\smallskip
\noindent\hangindent=1em Einige Bemerkungen zur Vereinfachung der Rechnung für die geocentrischen Oerter der Planeten \mbox{.\ .\ .\ .\ .\ .\ .\ .\ .\ .\ .\ .\ .\ .\ .\ .\ .\ .\ .} \hfill\mbox{1804 Mai}\quad\mbox{--- 94}\par\smallskip
\noindent\hangindent=1em Ueber die Grenzen der geocentrischen Oerter der Planeten \mbox{.\ .\ .\ .\ .\ .} \hfill\mbox{1804 Aug.}\quad\mbox{. --- 106}\par\smallskip
\noindent\hangindent=1em Der Zodiacus der Juno \mbox{.\ .\ .\ .\ .\ .\ .\ .\ .\ .\ .\ .\ .\ .\ .\ .\ .\ .\ .} \hfill\mbox{1805 März}\quad\mbox{--- 119}\par\smallskip
\endgroup
